\section{Hardware-Specific Optimizations}
\label{sec:hardware}

\subsection{AMD Ryzen 7 PRO 6850H Analysis}

The Ryzen 7 PRO 6850H is a mobile processor built on AMD's Zen 3+ architecture. Understanding its characteristics is essential for effective optimization.

\subsubsection{Microarchitecture}

Key architectural features relevant to LLM inference:

\begin{table}[h]
\centering
\caption{Ryzen 7 PRO 6850H Key Specifications}
\label{tab:ryzen_specs}
\begin{tabular}{ll}
\toprule
\textbf{Feature} & \textbf{Specification} \\
\midrule
Architecture & Zen 3+ (Rembrandt) \\
Cores / Threads & 8 / 16 \\
Base Clock & 3.2 GHz \\
Max Boost & 4.7 GHz \\
L1 Cache (per core) & 32 KB + 32 KB (I$+$D) \\
L2 Cache (per core) & 512 KB \\
L3 Cache (shared) & 16 MB \\
Memory Support & DDR5-4800, LPDDR5-6400 \\
TDP & 35-45W configurable \\
SIMD Support & AVX2 (256-bit) \\
\bottomrule
\end{tabular}
\end{table}

The Zen 3+ architecture provides significant improvements over Zen 2 for our workload:

\begin{itemize}
    \item Improved branch prediction reduces misprediction penalties in attention computation
    \item Unified 8-way L3 cache reduces latency for data sharing between cores
    \item AVX2 support enables 256-bit vector operations, doubling throughput over AVX
\end{itemize}

\subsubsection{Memory Subsystem}

The integrated memory controller on the 6850H supports DDR5-4800, providing theoretical bandwidth:

\begin{equation}
B_{max} = 4800 \text{ MT/s} \times 64 \text{ bits} \times 2 \text{ channels} / 8 = 76.8 \text{ GB/s}
\end{equation}

However, practical sustained bandwidth is typically 70-75\% of theoretical due to protocol overhead, refresh cycles, and bank conflicts. We estimate effective bandwidth at approximately 50-55 GB/s.

For FP16 matrix-vector multiplication with $d_{model} = 3072$:

\begin{itemize}
    \item Weight matrix size: $3072 \times 3072 \times 2 \text{ bytes} = 18$ MB
    \item Input vector size: $3072 \times 2 \text{ bytes} = 6$ KB
    \item Output vector size: $3072 \times 2 \text{ bytes} = 6$ KB
    \item Total memory traffic per operation: $\approx 18$ MB
\end{itemize}

With 50 GB/s bandwidth, this suggests a theoretical maximum of approximately 2,778 operations/second, or for 32 layers, about 87 layer computations/second. At 3072 dimensions with FP16 (2 bytes), this translates to roughly 15-20 milliseconds per token for the matrix operations alone, matching our observed 15-25 tok/s generation rate.

\subsubsection{Cache Hierarchy Effects}

The cache hierarchy significantly impacts inference performance:

\begin{itemize}
    \item \textbf{L1 Cache:} 32KB instruction + 32KB data per core. Too small for weight matrices but effective for temporary activations.
    \item \textbf{L2 Cache:} 512KB per core. Can hold portions of weight matrices; tiling strategies can exploit this.
    \item \textbf{L3 Cache:} 16MB shared. Can hold approximately 896 rows of a 3072x3072 FP16 matrix. This is substantial for batched inference but limited for single-sequence generation.
\end{itemize}

\subsection{Platform-Specific Configuration}

The M7Config class encapsulates hardware-specific optimizations for the GMKtech M7 platform:

\begin{table}[h]
\centering
\caption{M7Config Hardware-Specific Settings}
\label{tab:m7config_settings}
\begin{tabular}{lll}
\toprule
\textbf{Parameter} & \textbf{Value} & \textbf{Rationale} \\
\midrule
n\_threads & 8 & Physical core count (no hyperthreading) \\
n\_batch & 2048 & Tuned for DDR5 bandwidth \\
n\_ubatch & 512 & L2 cache-friendly micro-batch \\
max\_memory\_gb & 24.0 & Available after iGPU allocation \\
reserve\_memory\_gb & 4.0 & OS and system overhead \\
cache\_type\_k & turbo3 & Key cache compression \\
cache\_type\_v & turbo3 & Value cache compression \\
flash\_attn & True & AMD-optimized attention \\
use\_mmap & True & Demand paging for weights \\
use\_mlock & False & WSL2 limitation (not supported) \\
\bottomrule
\end{tabular}
\end{table}

These defaults are derived from extensive profiling on the target hardware, balancing throughput, latency, and memory efficiency.

\subsection{WSL2 Considerations}

Running under Windows Subsystem for Linux 2 introduces specific constraints:

\subsubsection{Memory Management}

WSL2 uses a virtual machine with dynamic memory allocation. Key implications:

\begin{itemize}
    \item \texttt{mlock} is not fully supported, preventing memory locking
    \item Memory ballooning can introduce latency during allocation
    \item The VMM adds overhead for page table management
\end{itemize}

We disable \texttt{mlock} and use memory-mapped files (mmap) for model weights, allowing the operating system to manage paging.

\subsubsection{File System Performance}

WSL2 provides improved file system performance over WSL1, but path translation and 9P protocol overhead can still impact model loading. We recommend:

\begin{itemize}
    \item Storing models on the Linux filesystem (\texttt{/home/user/}) rather than Windows mounts (\texttt{/mnt/c/})
    \item Using asynchronous loading to hide I/O latency
    \item Enabling \texttt{use\_mmap} for demand paging of model weights
\end{itemize}

\subsection{AMD-Specific Optimizations}

\subsubsection{ZenDNN Integration}

While ZenDNN (AMD's deep learning library) primarily targets PyTorch and TensorFlow, llama.cpp can leverage some Zen architecture optimizations:

\begin{itemize}
    \item Optimized GEMM kernels for Zen 3+ matrix operations
    \item Improved memory access patterns for the L3 cache architecture
    \item Better thread scheduling across CCX complexes
\end{itemize}

\subsubsection{AVX2 Kernel Selection}

llama.cpp provides multiple GEMM implementations. For Ryzen processors:

\begin{itemize}
    \item Prefer \texttt{GGML\_CPU\_CAP\_AVX2} kernels
    \item Avoid AVX-512 code paths (not supported, emulation is slow)
    \item Enable FMA3 for fused multiply-add operations
\end{itemize}

\subsection{Power and Thermal Management}

The 6850H supports configurable TDP (cTDP) from 35W to 54W. For sustained inference:

\begin{itemize}
    \item Lower TDP settings reduce thermal throttling risk
    \item However, sustained boost clocks require adequate cooling
    \item The GMKtech M7 compact form factor limits sustained boost
\end{itemize}

We configure for sustained operation at base clock (3.2 GHz) rather than peak boost (4.7 GHz), providing predictable performance without thermal concerns.

\subsection{Unified Memory Architecture}

The integrated Radeon 680M graphics shares system memory, with 4GB allocated as VRAM:

\begin{itemize}
    \item Reduces available system RAM from 32GB to 28GB
    \item Eliminates PCIe data copy overhead for GPU inference
    \item Enables zero-copy for vision preprocessing
\end{itemize}

While we target CPU inference, the unified memory enables efficient multimodal processing where vision features can be computed on iGPU while language processing occurs on CPU.

\subsection{Performance Profiling}

We use AMD uProf and Linux perf for hardware-level profiling:

\begin{itemize}
    \item L3 cache miss rate: Target $<$5\% for weight matrix access
    \item AVX2 utilization: Target $>$90\% of theoretical FLOPS
    \item Memory bandwidth: Monitor for saturation
    \item Branch misprediction: Critical for attention softmax
\end{itemize}

Initial profiling revealed attention softmax as a bottleneck due to its sequential dependency. Flash Attention implementation addresses this by fusing operations and reducing HBM accesses.

\subsection{Optimal Thread Configuration}

Thread count significantly impacts performance due to synchronization overhead and cache effects:

\begin{table}[h]
\centering
\caption{Thread Count Performance Comparison}
\label{tab:thread_comparison}
\begin{tabular}{cccc}
\toprule
\textbf{Threads} & \textbf{Throughput (tok/s)} & \textbf{Latency (ms)} & \textbf{Efficiency} \\
\midrule
4 & 12.5 & 80 & 100\% \\
8 & 22.3 & 45 & 89\% \\
12 & 24.1 & 41 & 64\% \\
16 & 23.8 & 42 & 50\% \\
\bottomrule
\end{tabular}
\end{table}

Based on these results, we default to 8 threads (physical cores), achieving near-optimal throughput without the efficiency loss of hyperthreading.

\subsection{Batch Size Tuning}

The batch size (\texttt{n\_batch}) affects both throughput and latency:

\begin{itemize}
    \item Larger batches improve matrix multiplication efficiency
    \item However, they increase memory consumption
    \item For interactive generation (batch=1), smaller batches reduce latency
\end{itemize}

We configure \texttt{n\_batch=2048} as a balance, enabling efficient prompt processing while maintaining reasonable memory footprint.
